\documentclass[11pt,a4paper]{article}
\usepackage[utf8]{inputenc}
\usepackage[T1]{fontenc}
\usepackage{geometry}
\geometry{top=1in, bottom=1in, left=1in, right=1in}
\usepackage{hyperref}
\usepackage{booktabs}
\usepackage{amsmath}
\usepackage{listings}
\usepackage{xcolor}
\usepackage{graphicx}
\usepackage{array}
\usepackage{fancyhdr}
\usepackage{titlesec}
\usepackage{longtable}

\hypersetup{
    colorlinks=true,
    linkcolor=purple,
    filecolor=magenta,      
    urlcolor=cyan,
    pdftitle=Genesis Design System (GDS),
    pdfpagemode=FullScreen,
}

\definecolor{codegray}{rgb}{0.96,0.96,0.96}
\definecolor{codeborder}{rgb}{0.85,0.85,0.85}
\definecolor{commentgreen}{rgb}{0.12,0.55,0.22}
\definecolor{keywordblue}{rgb}{0.1,0.2,0.8}

\lstset{
    backgroundcolor=\color{codegray},
    basicstyle=\ttfamily\small,
    breakatwhitespace=false,
    breaklines=true,
    captionpos=b,
    commentstyle=\color{commentgreen},
    frame=single,
    rulecolor=\color{codeborder},
    keepspaces=true,
    keywordstyle=\color{keywordblue},
    language=Python,
    showspaces=false,
    showstringspaces=false,
    showtabs=false,
    tabsize=4
}

\pagestyle{fancy}
\fancyhf{}
\rhead{\small Genesis Design System (GDS)}
\lhead{\small Project Genesis}
\cfoot{\thepage}

\title{\textbf{Genesis Design System (GDS)}}
\author{Project Genesis Research Repository}
\date{\small Generated: July 2026}

\begin{document}

\maketitle
\thispagestyle{empty}

\newpage
\section*{Genesis Design System (GDS)}
\subsection*{Version 1.0  Visual Language, Tokens \& Components}

\textbackslash{}textbf{Project:} Project Genesis  Agentic Artificial Life \& Evolutionary Biology Simulator
\textbackslash{}textbf{Document Role:} Visual design system  colors, typography, animations, components
\textbackslash{}textbf{Companion Documents:} GPS (Product)  GES (Engineering)  GCR (Content \& Research)
\textbackslash{}textbf{Status:} Living document  extend as new components are built

\noindent\rule{\textwidth}{0.4pt}

\textgreater{} "Design is not decoration. It is the architecture of experience."

\noindent\rule{\textwidth}{0.4pt}

\subsection*{Part I  Design Philosophy}

\subsubsection*{1.1 The Visual Identity of Genesis}

Genesis exists at the intersection of \textbackslash{}textbf{science and wilderness}.

It is a research tool  precise, structured, evidence-based.
It is also a window into a world  dramatic, alive, unknowable.

The design must hold both truths simultaneously.

The result is a visual language built from three sources:

\begin{enumerate}
  \item \textbackslash{}textbf{Space and deep time.} Dark backgrounds. Sparse light. Long durations. The feeling of watching something enormous from a distance.
\end{enumerate}

\begin{enumerate}
  \item \textbackslash{}textbf{Bioluminescence.} Teal accents. Green-to-teal gradients. The visual register of life in dark environments  organisms that glow because they have to.
\end{enumerate}

\begin{enumerate}
  \item \textbackslash{}textbf{Scientific instruments.} Monospace data. Grid-aligned layouts. Precise labeling. The aesthetic of a spectrometer readout or a seismograph.
\end{enumerate}

These three sources produce a visual language that is simultaneously beautiful and rigorous.

\subsubsection*{1.2 One Mode Only}

Genesis is \textbackslash{}textbf{dark mode only}.

Not because dark mode is fashionable.

Because you do not look at stars in daylight. Because laboratory instruments are dark. Because the world in Genesis is rendered on a dark canvas. The light elements  the teal rivers, the amber colony, the red death events  only sing against darkness.

There is no light mode. There is no toggle.

\noindent\rule{\textwidth}{0.4pt}

\subsection*{Part II  Color System}

\subsubsection*{2.1 Background Stack}

\begin{lstlisting}[language=css]
/* These four layers create the depth hierarchy of the interface */

--bg-void:      hsl(222, 30%, 4%);    /* #070A10 — absolute dark, page background */
--bg-primary:   hsl(222, 28%, 7%);    /* #0C1019 — primary surface */
--bg-secondary: hsl(222, 24%, 10%);   /* #141824 — slightly elevated cards */
--bg-surface:   hsl(222, 20%, 14%);   /* #1C2130 — elevated elements (inputs, etc.) */
--bg-elevated:  hsl(222, 18%, 18%);   /* #242C3D — highest elevation (modals) */
\end{lstlisting}

\textbackslash{}textbf{Rule:} Each surface level is exactly one stop above the previous. Never skip levels.

\subsubsection*{2.2 Accent  Bioluminescent Teal}

\begin{lstlisting}[language=css]
--accent-50:   hsl(172, 80%, 95%);   /* Near white teal — for text on teal bg */
--accent-100:  hsl(172, 75%, 88%);
--accent-200:  hsl(172, 75%, 75%);
--accent-300:  hsl(172, 75%, 62%);
--accent-400:  hsl(172, 80%, 52%);   /* Main accent — use for CTAs, active states */
--accent-500:  hsl(172, 80%, 45%);   /* Hover state */
--accent-600:  hsl(172, 70%, 36%);   /* Pressed state */
--accent-700:  hsl(172, 60%, 28%);   /* Borders, dim accents */
--accent-800:  hsl(172, 50%, 18%);   /* Subtle accent backgrounds */
--accent-900:  hsl(172, 40%, 10%);   /* Very subtle accent tint */

--accent-glow-sm:  0 0 12px hsl(172, 80%, 48%, 0.20);
--accent-glow-md:  0 0 24px hsl(172, 80%, 48%, 0.25);
--accent-glow-lg:  0 0 48px hsl(172, 80%, 48%, 0.30);
\end{lstlisting}

\subsubsection*{2.3 Colony Colors}

These match the simulation's colony color system exactly. Never change them  they must be consistent between the visualizer and the website.

\begin{lstlisting}[language=css]
--colony-alpha:       hsl(210, 85%, 58%);    /* Cobalt Blue */
--colony-alpha-dim:   hsl(210, 60%, 30%);
--colony-alpha-glow:  hsl(210, 85%, 58%, 0.20);

--colony-beta:        hsl(35,  88%, 58%);    /* Warm Amber */
--colony-beta-dim:    hsl(35,  60%, 30%);
--colony-beta-glow:   hsl(35,  88%, 58%, 0.20);

--colony-gamma:       hsl(280, 78%, 63%);    /* Violet */
--colony-gamma-dim:   hsl(280, 55%, 30%);
--colony-gamma-glow:  hsl(280, 78%, 63%, 0.20);

--colony-delta:       hsl(140, 68%, 52%);    /* Emerald Green */
--colony-delta-dim:   hsl(140, 50%, 26%);
--colony-delta-glow:  hsl(140, 68%, 52%, 0.20);
\end{lstlisting}

\subsubsection*{2.4 Semantic Event Colors}

\begin{lstlisting}[language=css]
--event-birth:          hsl(120, 58%, 52%);    /* Verdant green */
--event-death:          hsl(0,   68%, 52%);    /* Blood red */
--event-milestone:      hsl(45,  88%, 58%);    /* Gold */
--event-disaster:       hsl(18,  88%, 52%);    /* Volcanic orange */
--event-dispute:        hsl(300, 58%, 58%);    /* Electric violet */
--event-climate:        hsl(195, 75%, 52%);    /* Ice blue */
--event-extinction:     hsl(0,   50%, 38%);    /* Dark rust — muted for gravity */

/* Dot shadows for event markers in timeline */
--event-birth-shadow:   0 0 6px hsl(120, 58%, 52%, 0.5);
--event-death-shadow:   0 0 6px hsl(0,   68%, 52%, 0.5);
--event-milestone-shadow: 0 0 6px hsl(45, 88%, 58%, 0.5);
\end{lstlisting}

\subsubsection*{2.5 Biome Colors}

\begin{lstlisting}[language=css]
/* Match visualizer.html palette exactly */
--biome-ocean:          hsl(207, 75%, 28%);
--biome-glacier:        hsl(200, 35%, 82%);
--biome-tundra:         hsl(200, 22%, 60%);
--biome-taiga:          hsl(142, 38%, 30%);
--biome-temperate:      hsl(118, 48%, 36%);
--biome-grassland:      hsl(78,  58%, 42%);
--biome-desert:         hsl(36,  65%, 50%);
--biome-rainforest:     hsl(152, 68%, 26%);
--biome-lake:           hsl(202, 72%, 48%);
\end{lstlisting}

\subsubsection*{2.6 Text Colors}

\begin{lstlisting}[language=css]
--text-primary:    hsl(220, 18%, 92%);    /* Near white — main body */
--text-secondary:  hsl(220, 14%, 68%);    /* Muted — secondary info */
--text-tertiary:   hsl(220, 10%, 48%);    /* Very muted — labels, placeholders */
--text-accent:     hsl(172, 80%, 52%);    /* Teal — links, highlights */
--text-inverse:    hsl(222, 28%, 7%);     /* Dark — text on teal buttons */
--text-danger:     hsl(0,   70%, 60%);    /* Red — errors */
--text-warning:    hsl(38,  85%, 58%);    /* Amber — warnings */
--text-success:    hsl(142, 65%, 52%);    /* Green — success states */
\end{lstlisting}

\subsubsection*{2.7 Border Colors}

\begin{lstlisting}[language=css]
--border-subtle:   hsl(222, 20%, 18%);    /* Barely visible — structural dividers */
--border-default:  hsl(222, 20%, 24%);    /* Default card/input borders */
--border-hover:    hsl(222, 20%, 32%);    /* Hover state borders */
--border-accent:   hsl(172, 60%, 32%);    /* Teal borders (active states, focus) */
--border-danger:   hsl(0,   60%, 40%);    /* Error state borders */
\end{lstlisting}

\subsubsection*{2.8 Gradients}

\begin{lstlisting}[language=css]
/* Primary CTA gradient */
--gradient-cta: linear-gradient(135deg,
  hsl(172, 80%, 42%) 0%,
  hsl(192, 80%, 38%) 100%);

/* Card accent top border gradient (applied as border-top) */
--gradient-card-top: linear-gradient(90deg,
  hsl(172, 80%, 48%) 0%,
  hsl(192, 70%, 45%) 50%,
  transparent 100%);

/* World map temperature gradient */
--gradient-temperature: linear-gradient(90deg,
  hsl(230, 80%, 35%) 0%,    /* Cold — arctic blue */
  hsl(195, 70%, 50%) 25%,   /* Cool — cerulean */
  hsl(55,  70%, 52%) 60%,   /* Warm — amber */
  hsl(15,  80%, 50%) 100%); /* Hot — red-orange */

/* World map elevation gradient */
--gradient-elevation: linear-gradient(90deg,
  hsl(207, 75%, 20%) 0%,    /* Deep ocean */
  hsl(207, 60%, 35%) 28%,   /* Shallow water */
  hsl(118, 40%, 32%) 32%,   /* Lowland green */
  hsl(78,  40%, 40%) 50%,   /* Mid-elevation */
  hsl(35,  30%, 50%) 70%,   /* Highland */
  hsl(0,   0%,  75%) 100%); /* Snow cap */
\end{lstlisting}

\noindent\rule{\textwidth}{0.4pt}

\subsection*{Part III  Typography}

\subsubsection*{3.1 Font Stack}

\begin{lstlisting}[language=css]
@import url('https://fonts.googleapis.com/css2?family=Space+Grotesk:wght@300;400;500;600;700&family=Inter:wght@400;500;600&family=JetBrains+Mono:wght@400;500&display=swap');

--font-display:  'Space Grotesk', system-ui, -apple-system, sans-serif;
--font-body:     'Inter', system-ui, -apple-system, sans-serif;
--font-mono:     'JetBrains Mono', 'Fira Code', 'Cascadia Code', monospace;
\end{lstlisting}

\textbackslash{}textbf{Space Grotesk}  All headings, UI labels, navigation, CTA buttons. Its slightly geometric, wide letterform reads as scientific without being cold.

\textbackslash{}textbf{Inter}  All body copy, descriptions, abstracts. Maximally readable at small sizes.

\textbackslash{}textbf{JetBrains Mono}  All data values, IDs, tick counts, gene values, config params, code blocks. Never use a proportional font for numbers that align vertically.

\subsubsection*{3.2 Type Scale}

\begin{lstlisting}[language=css]
--text-2xs:  0.625rem;   /* 10px — data labels, legend text */
--text-xs:   0.75rem;    /* 12px — captions, secondary metadata */
--text-sm:   0.875rem;   /* 14px — secondary body, UI labels */
--text-base: 1rem;       /* 16px — primary body copy */
--text-lg:   1.125rem;   /* 18px — lead text, card descriptions */
--text-xl:   1.25rem;    /* 20px — card titles, section labels */
--text-2xl:  1.5rem;     /* 24px — page section headers */
--text-3xl:  1.875rem;   /* 30px — page titles */
--text-4xl:  2.25rem;    /* 36px — major section heroes */
--text-5xl:  3rem;       /* 48px — landing hero text */
--text-6xl:  3.75rem;    /* 60px — ultra-large display (sparingly) */
\end{lstlisting}

\subsubsection*{3.3 Font Weight Usage}

\begin{longtable}{lll}
\toprule
Use Case & Font & Weight \\
\midrule
\endhead
Page title & Space Grotesk & 700 \\
Section heading (h2) & Space Grotesk & 600 \\
Card title & Space Grotesk & 600 \\
Navigation links & Space Grotesk & 500 \\
Body copy & Inter & 400 \\
Secondary body & Inter & 400 \\
Lead paragraph & Inter & 400 (but text-lg) \\
Button label & Space Grotesk & 600 \\
Metric value & JetBrains Mono & 500 \\
Data label & Inter & 500 \\
Code/config & JetBrains Mono & 400 \\
Phase badge & Space Grotesk & 700 \\
\bottomrule
\end{longtable}

\subsubsection*{3.4 Line Heights}

\begin{lstlisting}[language=css]
--leading-none:     1;
--leading-tight:    1.2;   /* Headlines */
--leading-snug:     1.35;  /* Card titles */
--leading-normal:   1.5;   /* UI elements */
--leading-relaxed:  1.7;   /* Body copy */
--leading-loose:    2.0;   /* Spacious reading (monograph) */
\end{lstlisting}

\subsubsection*{3.5 Letter Spacing}

\begin{lstlisting}[language=css]
--tracking-tighter: -0.03em;   /* Large display headings */
--tracking-tight:   -0.01em;   /* Regular headings */
--tracking-normal:  0;
--tracking-wide:    0.04em;    /* UI labels, phase badges */
--tracking-wider:   0.08em;    /* All-caps micro-labels */
--tracking-widest:  0.14em;    /* Section category labels (e.g. "EXPERIMENT") */
\end{lstlisting}

\subsubsection*{3.6 Prose Styles (Research Documents)}

For the monograph, architecture doc, lab notebook, and thesis pages:

\begin{lstlisting}[language=css]
.prose {
  font-family: var(--font-body);
  font-size: var(--text-base);
  line-height: var(--leading-loose);
  color: var(--text-primary);
  max-width: var(--content-width);   /* 800px */

  h1 { font: 700 var(--text-4xl)/var(--leading-tight) var(--font-display); margin-bottom: 2rem; }
  h2 { font: 600 var(--text-2xl)/var(--leading-tight) var(--font-display); margin-top: 3rem; margin-bottom: 1rem; }
  h3 { font: 600 var(--text-xl)/var(--leading-snug) var(--font-display); margin-top: 2rem; }

  p    { margin-bottom: 1.5rem; }
  code { font-family: var(--font-mono); font-size: 0.9em; background: var(--bg-elevated); padding: 0.1em 0.4em; border-radius: var(--radius-sm); }
  pre  { background: var(--bg-elevated); border: 1px solid var(--border-default); border-radius: var(--radius-md); padding: 1.5rem; overflow-x: auto; }

  blockquote {
    border-left: 3px solid var(--accent-700);
    padding-left: 1.5rem;
    color: var(--text-secondary);
    font-style: italic;
  }

  table { width: 100%; border-collapse: collapse; }
  th, td { padding: 0.75rem 1rem; border-bottom: 1px solid var(--border-subtle); text-align: left; }
  th { color: var(--text-secondary); font: 500 var(--text-sm)/1 var(--font-body); letter-spacing: var(--tracking-wider); text-transform: uppercase; }
}
\end{lstlisting}

\noindent\rule{\textwidth}{0.4pt}

\subsection*{Part IV  Spacing \& Layout}

\subsubsection*{4.1 Spacing Scale (4px base unit)}

\begin{lstlisting}[language=css]
--space-0:   0;
--space-px:  1px;
--space-0-5: 0.125rem;  /* 2px */
--space-1:   0.25rem;   /* 4px */
--space-2:   0.5rem;    /* 8px */
--space-3:   0.75rem;   /* 12px */
--space-4:   1rem;      /* 16px */
--space-5:   1.25rem;   /* 20px */
--space-6:   1.5rem;    /* 24px */
--space-7:   1.75rem;   /* 28px */
--space-8:   2rem;      /* 32px */
--space-10:  2.5rem;    /* 40px */
--space-12:  3rem;      /* 48px */
--space-14:  3.5rem;    /* 56px */
--space-16:  4rem;      /* 64px */
--space-20:  5rem;      /* 80px */
--space-24:  6rem;      /* 96px */
--space-32:  8rem;      /* 128px */
\end{lstlisting}

\subsubsection*{4.2 Border Radius}

\begin{lstlisting}[language=css]
--radius-none: 0;
--radius-sm:   3px;
--radius-md:   6px;
--radius-lg:   10px;
--radius-xl:   14px;
--radius-2xl:  20px;
--radius-3xl:  28px;
--radius-full: 9999px;   /* Pills, fully circular */
\end{lstlisting}

\subsubsection*{4.3 Layout Grid}

\begin{lstlisting}[language=css]
--max-width:       1280px;   /* Page container */
--content-width:   800px;    /* Prose content */
--sidebar-width:   240px;    /* Experiment detail sidebar */
--grid-gap:        1.5rem;   /* 24px — standard grid gap */
--grid-gap-sm:     1rem;     /* 16px — tight grids */
--grid-gap-lg:     2rem;     /* 32px — loose grids */
\end{lstlisting}

\subsubsection*{4.4 Responsive Breakpoints}

\begin{lstlisting}[language=css]
/* Mobile-first approach */
--bp-xs:  360px;    /* Small phones */
--bp-sm:  640px;    /* Large phones, small tablets */
--bp-md:  768px;    /* Tablets */
--bp-lg:  1024px;   /* Small laptops */
--bp-xl:  1280px;   /* Desktops */
--bp-2xl: 1536px;   /* Large screens */
\end{lstlisting}

\subsubsection*{4.5 Mobile Rules}

\begin{itemize}
  \item Sidebar becomes bottom tab bar (horizontal scroll) on \textless{} \textbackslash{}texttt{--bp-lg}
  \item Replay canvas fills full viewport width on mobile
  \item World map uses touch pan/pinch instead of mouse
  \item Timeline events stack full-width
  \item Metric cards: 2 columns on mobile, 4 on desktop
  \item Archive grid: 1 column on mobile, 2 on tablet, 3 on desktop
\end{itemize}

\noindent\rule{\textwidth}{0.4pt}

\subsection*{Part V  Effects \& Surfaces}

\subsubsection*{5.1 Glass Morphism}

Used on: cards, sidebars, modals, tooltip panels, the agent inspector.

\begin{lstlisting}[language=css]
.glass {
  background: hsl(222, 24%, 10%, 0.75);
  backdrop-filter: blur(20px) saturate(160%);
  -webkit-backdrop-filter: blur(20px) saturate(160%);
  border: 1px solid hsl(222, 20%, 22%);
  border-top-color: hsl(222, 20%, 28%);   /* Subtle top highlight */
}

.glass-strong {
  background: hsl(222, 24%, 8%, 0.88);
  backdrop-filter: blur(32px) saturate(180%);
}
\end{lstlisting}

\subsubsection*{5.2 Shadows}

\begin{lstlisting}[language=css]
--shadow-none: none;
--shadow-sm:   0 1px 2px hsl(222, 40%, 2%, 0.40),
               0 1px 4px hsl(222, 40%, 2%, 0.20);
--shadow-md:   0 4px 8px  hsl(222, 40%, 2%, 0.45),
               0 2px 16px hsl(222, 40%, 2%, 0.25);
--shadow-lg:   0 8px 16px hsl(222, 40%, 2%, 0.50),
               0 4px 32px hsl(222, 40%, 2%, 0.30);
--shadow-xl:   0 16px 32px hsl(222, 40%, 2%, 0.55),
               0 8px 64px  hsl(222, 40%, 2%, 0.35);

/* Accent glow shadows */
--shadow-accent-sm: 0 0 12px hsl(172, 80%, 48%, 0.18);
--shadow-accent-md: 0 0 24px hsl(172, 80%, 48%, 0.22);
--shadow-accent-lg: 0 0 48px hsl(172, 80%, 48%, 0.28);
\end{lstlisting}

\subsubsection*{5.3 Elevation System}

\begin{longtable}{lllll}
\toprule
Elevation & Surface & Shadow & Blur & Use Case \\
\midrule
\endhead
0 & \textbackslash{}texttt{--bg-void} & none & none & Page background \\
1 & \textbackslash{}texttt{--bg-primary} & \textbackslash{}texttt{shadow-sm} & none & Main content area \\
2 & \textbackslash{}texttt{--bg-secondary} & \textbackslash{}texttt{shadow-md} & none & Cards, panels \\
3 & \textbackslash{}texttt{--bg-surface} & \textbackslash{}texttt{shadow-lg} & glass & Dropdowns, tooltips \\
4 & \textbackslash{}texttt{--bg-elevated} & \textbackslash{}texttt{shadow-xl} & glass-strong & Modals, overlays \\
5 & Accent & \textbackslash{}texttt{shadow-accent-md} & none & Featured, active, selected \\
\bottomrule
\end{longtable}

\noindent\rule{\textwidth}{0.4pt}

\subsection*{Part VI  Animation System}

\subsubsection*{6.1 Easing Functions}

\begin{lstlisting}[language=css]
--ease-linear:      linear;
--ease-in:          cubic-bezier(0.4, 0, 1, 1);
--ease-out:         cubic-bezier(0, 0, 0.2, 1);
--ease-in-out:      cubic-bezier(0.4, 0, 0.2, 1);

/* Genesis-specific easings */
--ease-out-quart:   cubic-bezier(0.25, 1, 0.5, 1);        /* Smooth, confident */
--ease-spring:      cubic-bezier(0.34, 1.56, 0.64, 1);    /* Springy — microinteractions */
--ease-overshoot:   cubic-bezier(0.68, -0.55, 0.27, 1.55);/* Playful overshoot */
--ease-slow-out:    cubic-bezier(0.16, 1, 0.3, 1);        /* Dramatic reveal */
\end{lstlisting}

\subsubsection*{6.2 Duration Scale}

\begin{lstlisting}[language=css]
--duration-instant:  50ms;    /* State flashes, immediate feedback */
--duration-fast:     100ms;   /* Hover states */
--duration-normal:   200ms;   /* Most transitions */
--duration-moderate: 300ms;   /* Panel slides, dropdown opens */
--duration-slow:     500ms;   /* Page section reveals */
--duration-slower:   700ms;   /* Hero elements, major transitions */
--duration-page:     1000ms;  /* Full-page transitions */
\end{lstlisting}

\subsubsection*{6.3 Micro-interaction Specifications}

Every interactive element is specified here. No guessing.

\textbackslash{}textbf{Button  primary (teal)}
\begin{lstlisting}[language=Python]
rest:    bg=accent-400, shadow=none, scale=1.0
hover:   bg=accent-300, shadow=accent-sm, scale=1.02  [150ms ease-out-quart]
active:  bg=accent-600, shadow=none,     scale=0.97   [100ms ease-in]
focus:   ring: 2px accent-400 offset 2px
\end{lstlisting}

\textbackslash{}textbf{Button  ghost}
\begin{lstlisting}[language=Python]
rest:    bg=transparent, border=border-default, color=text-secondary
hover:   bg=bg-surface, border=border-hover, color=text-primary  [150ms ease-out]
active:  bg=bg-elevated  [100ms ease-in]
\end{lstlisting}

\textbackslash{}textbf{Experiment Card}
\begin{lstlisting}[language=Python]
rest:    bg=bg-secondary, border=border-subtle, shadow=shadow-sm, translateY=0
hover:   bg=bg-surface, border=border-hover, shadow=shadow-md, translateY=-3px  [200ms ease-out-quart]
         thumbnail: scale=1.04 (inner only, parent overflow:hidden)              [400ms ease-out-quart]
         world preset badge: color transitions to accent                          [150ms ease-out]
active:  translateY=-1px, shadow=shadow-sm  [100ms ease-in]
featured: border-top: 2px solid accent-400 (always on)
\end{lstlisting}

\textbackslash{}textbf{Timeline Event Row}
\begin{lstlisting}[language=Python]
rest:    opacity=1, pl=var(--space-4)
hover:   bg=hsl(222,20%,14%,0.5), pl=var(--space-5) (left nudge)  [150ms ease-out]
         event dot: scale=1.3, glow appears  [200ms ease-spring]
active:  bg=bg-surface
\end{lstlisting}

\textbackslash{}textbf{Tab (sidebar navigation)}
\begin{lstlisting}[language=Python]
rest:    color=text-secondary, bg=transparent
hover:   color=text-primary, bg=hsl(222,20%,14%)  [100ms ease-out]
active:  color=accent-400, bg=accent-900
         left border: 2px solid accent-400
\end{lstlisting}

\textbackslash{}textbf{Layer Toggle Chip (world map)}
\begin{lstlisting}[language=Python]
rest:    bg=bg-elevated, color=text-secondary, border=border-default
hover:   bg=bg-surface, color=text-primary  [150ms ease-out]
active:  bg=accent-900, color=accent-300, border=accent-700
toggle:  layer fades in/out on canvas: 200ms ease-out
\end{lstlisting}

\textbackslash{}textbf{World Map Agent Dot}
\begin{lstlisting}[language=Python]
render:  radius=3px baseline; filled with colony color
         opacity = agent.health / 100  (full health = fully opaque)
move:    lerp between tick positions (linear, per-frame)
hover:   radius expands to 6px  [150ms ease-spring]
         inspector panel slides in from right  [300ms ease-out-quart]
death:   dot flashes white (50ms), then red (200ms), then fades to 0 opacity (300ms)
birth:   dot appears at parent midpoint, scale 0→1  [300ms ease-spring]
\end{lstlisting}

\textbackslash{}textbf{Metric Counter (on viewport enter)}
\begin{lstlisting}[language=Python]
trigger: IntersectionObserver, 0.1 threshold
effect:  value counts from 0 to final value
easing:  ease-out-quart applied to the count interpolation
duration: 800ms
\end{lstlisting}

\textbackslash{}textbf{Search Bar}
\begin{lstlisting}[language=Python]
rest:    width=200px, bg=bg-secondary, border=border-default
focus:   width=320px (animated), border=border-accent, shadow=accent-sm  [250ms ease-out-quart]
results: slide in from below, opacity 0→1  [200ms ease-out]
\end{lstlisting}

\textbackslash{}textbf{Modal}
\begin{lstlisting}[language=Python]
open:    backdrop fades in (bg:hsl(222,40%,2%,0.7)) [250ms ease-out]
         panel: scale 0.95→1.0, opacity 0→1  [300ms ease-out-quart]
close:   reverse, both  [200ms ease-in]
\end{lstlisting}

\textbackslash{}textbf{Toast Notification}
\begin{lstlisting}[language=Python]
enter:   slide up from bottom-right (translateY: 80px→0), opacity 0→1  [300ms ease-spring]
auto-exit: after 4000ms, slide down + fade  [200ms ease-in]
hover:   auto-exit timer pauses
\end{lstlisting}

\textbackslash{}textbf{Replay Controls}
\begin{lstlisting}[language=Python]
play/pause: icon morphs with SVG path animation  [200ms ease-in-out]
speed:  click cycles through 0.5×→1×→2×→4× with badge update  [100ms]
tick counter: monospace, updates each rendered frame
scrubber: drag to seek — canvas re-renders at target tick on mouseup
\end{lstlisting}

\subsubsection*{6.4 Page-Level Transitions}

\begin{lstlisting}[language=Python]
Route change:
  outgoing: opacity 1→0, translateY 0→(-8px)  [200ms ease-in]
  incoming: opacity 0→1, translateY (8px)→0   [300ms ease-out]
  total: 500ms (overlap at 200ms)

Section reveals (scroll):
  trigger: IntersectionObserver, threshold 0.15
  effect: opacity 0→1, translateY 20px→0  [500ms ease-slow-out]
  stagger: 80ms between sibling elements
\end{lstlisting}

\subsubsection*{6.5 Canvas Animation  Replay Renderer}

\begin{lstlisting}[language=javascript]
// Per-frame rendering contract
function renderFrame(tick, agents, worldLayers, config) {
  // 1. Clear canvas
  ctx.clearRect(0, 0, canvas.width, canvas.height);

  // 2. Render base layer (biome map — pre-rendered ImageData, don't recompute)
  ctx.putImageData(biomeCacheImageData, 0, 0);

  // 3. Render optional overlay layers (rivers, resources, habitability)
  if (layers.rivers)      renderRiverOverlay(ctx, worldState);
  if (layers.habitability) renderHeatmap(ctx, worldState.habitability, 0.4);

  // 4. Render death density (persistent death heatmap)
  if (layers.deaths) renderHeatmap(ctx, worldState.death_density, 0.35);

  // 5. Render agents
  for (const agent of agents) {
    const lerped = lerpPosition(agent.prevPos, agent.pos, frameProgress);
    const alpha  = Math.max(0.2, agent.health / 100);
    const radius = agent.hovered ? 6 : 3;
    const color  = COLONY_COLORS[agent.colony_id];

    ctx.beginPath();
    ctx.arc(lerped.x, lerped.y, radius, 0, Math.PI * 2);
    ctx.fillStyle = `rgba(${color.r}, ${color.g}, ${color.b}, ${alpha})`;
    ctx.fill();

    // Death flash animation
    if (agent.deathFlash > 0) {
      ctx.fillStyle = `rgba(255, 255, 255, ${agent.deathFlash})`;
      ctx.fill();
      agent.deathFlash = Math.max(0, agent.deathFlash - 0.08);
    }
  }

  // 6. Render hovered agent highlight ring
  if (hoveredAgent) {
    ctx.strokeStyle = `rgba(255, 255, 255, 0.8)`;
    ctx.lineWidth = 1.5;
    ctx.beginPath();
    ctx.arc(hoveredAgent.x, hoveredAgent.y, 8, 0, Math.PI * 2);
    ctx.stroke();
  }
}
\end{lstlisting}

\noindent\rule{\textwidth}{0.4pt}

\subsection*{Part VII  Icon Vocabulary}

Genesis uses a minimal, consistent icon set. No icon packs. Custom SVGs where needed.

\subsubsection*{7.1 Event Type Icons}

\begin{longtable}{lll}
\toprule
Event & Icon & Color \\
\midrule
\endhead
Birth & \textbackslash{}texttt{} (expanding ring) & \textbackslash{}texttt{--event-birth} \\
Death & \textbackslash{}texttt{} (cross) & \textbackslash{}texttt{--event-death} \\
Milestone & \textbackslash{}texttt{} (diamond) & \textbackslash{}texttt{--event-milestone} \\
Disaster & \textbackslash{}texttt{} (triangle, storm) & \textbackslash{}texttt{--event-disaster} \\
Dispute & \textbackslash{}texttt{} (crossed swords, custom SVG) & \textbackslash{}texttt{--event-dispute} \\
Climate Epoch & \textbackslash{}texttt{\textbackslash{}textasciitilde{}} (wave) & \textbackslash{}texttt{--event-climate} \\
Extinction & \textbackslash{}texttt{} (solid square, black) & \textbackslash{}texttt{--event-extinction} \\
\bottomrule
\end{longtable}

\subsubsection*{7.2 Navigation Icons}

\begin{longtable}{ll}
\toprule
Page & Icon \\
\midrule
\endhead
Archive & Grid (22) \\
Observatory & Circle with dot (eye symbol) \\
Replay & Triangle (play) \\
Chronicle & Vertical line with ticks \\
World & Globe outline \\
Lineage & Tree branch \\
Genome & Helix (simplified) \\
Conditions & Sliders \\
Research & Book open \\
Timeline & Horizontal line with dots \\
Origin & Sparkle / star \\
\bottomrule
\end{longtable}

\subsubsection*{7.3 Icon Rules}

\begin{itemize}
  \item Size: 16px for inline, 20px for navigation, 24px for feature icons
  \item Stroke: 1.5px, round linecap, round linejoin
  \item Never fill + stroke together (pick one)
  \item Never use icon fonts  always SVG
\end{itemize}

\noindent\rule{\textwidth}{0.4pt}

\subsection*{Part VIII  Component Visual Specifications}

\subsubsection*{8.1 Navbar}

\begin{lstlisting}[language=Python]
Height: 60px
Position: fixed top, full width
Background: hsl(222, 28%, 6%, 0.85) + blur(20px)
Border-bottom: 1px solid border-subtle (always present, not on-scroll)

Left:
  Genesis wordmark — Space Grotesk 700, 18px, accent-400
  Phase badge — pill, bg:accent-900, color:accent-400, font:mono 11px

Center:
  Nav links — Space Grotesk 500, 14px, text-secondary
  Hover: color→text-primary + underline (2px accent-700 bottom)
  Active: color→accent-400 + underline (2px accent-400)
  Gap between links: 32px

Right:
  GitHub icon — 20px, text-tertiary → text-primary on hover
  "Command" button (admin only) — teal pill button, 500 font, 14px
\end{lstlisting}

\subsubsection*{8.2 Experiment/Civilization Card}

\begin{lstlisting}[language=Python]
Width: flexible (grid column)
Border-radius: radius-xl (14px)
Border: 1px solid border-subtle
Background: bg-secondary

Thumbnail area:
  Height: 140px
  overflow: hidden
  Object: world.png, cover fit
  Overlay: linear-gradient(to bottom, transparent 60%, bg-secondary 100%)

Content area:
  Padding: space-5 (20px)

  ID row:
    font: mono 11px, text-tertiary, tracking-wide
    "EXP-20260627-XY81"

  Title:
    font: Space Grotesk 600, text-xl, text-primary
    max 2 lines, line-clamp

  Tags row:
    Pills: bg-elevated, border-subtle, font:Inter 12px, text-secondary
    Gap: space-2

  Stats row (mono):
    "18 agents · 7,970 ticks (22 yr)"
    font: JetBrains Mono 12px, text-tertiary

  Meta row:
    "Published Jul 13, 2026"
    font: Inter 12px, text-tertiary

Featured variant:
  border-top: 2px solid gradient-card-top
  shadow: shadow-accent-sm
\end{lstlisting}

\subsubsection*{8.3 Metric Card}

\begin{lstlisting}[language=Python]
Padding: space-6
Border-radius: radius-lg
Background: bg-secondary
Border: 1px solid border-subtle

Label:
  font: Inter 500, text-xs, text-tertiary, tracking-wider, uppercase

Value:
  font: JetBrains Mono 500, text-3xl, text-primary
  margin-top: space-2

Unit (if any):
  font: Inter 400, text-sm, text-secondary, inline after value

Delta (optional):
  font: Inter 500, text-sm
  color: green if positive, red if negative
  "↑ +12% vs average"
  margin-top: space-1

Accent variant:
  border-left: 3px solid accent-700

Warning variant:
  border-left: 3px solid event-disaster

Critical variant:
  border-left: 3px solid event-death
\end{lstlisting}

\subsubsection*{8.4 Phase Badge (Timeline)}

\begin{lstlisting}[language=Python]
Display: inline-flex, align-center
Padding: space-1 space-3
Border-radius: radius-full
font: Space Grotesk 700, text-xs, tracking-wider, uppercase

Phase 0-1: bg=hsl(222,20%,20%), color=text-secondary
Phase 2-4: bg=hsl(195,30%,18%), color=hsl(195,70%,65%)
Phase 5-7: bg=hsl(172,30%,14%), color=accent-300
Phase 8-9: bg=accent-900, color=accent-400
Phase 10+: bg=hsl(280,30%,18%), color=hsl(280,60%,70%)
\end{lstlisting}

\subsubsection*{8.5 Code/Config Viewer}

\begin{lstlisting}[language=Python]
Background: bg-elevated
Border: 1px solid border-default
Border-radius: radius-md
Padding: space-6

Header bar:
  filename (mono, text-sm, text-secondary)
  Copy button (right-aligned)

Content:
  font: JetBrains Mono 400, text-sm, leading-relaxed
  JSON syntax highlighting:
    keys:    accent-300
    strings: hsl(45, 80%, 65%)    (amber)
    numbers: hsl(195, 75%, 65%)   (light blue)
    booleans: event-birth / event-death
    null:    text-tertiary
\end{lstlisting}

\noindent\rule{\textwidth}{0.4pt}

\subsection*{Part IX  World Map Rendering Specification}

\subsubsection*{9.1 Canvas Setup}

\begin{lstlisting}[language=Python]
Resolution: 1:1 cell-to-pixel mapping (512×512 for a 512-wide world)
Scaling: CSS transforms for zoom (not canvas resize — avoid re-render on zoom)
Pan: transform-origin center, translate on drag
Cursor: crosshair default, grab on pan
\end{lstlisting}

\subsubsection*{9.2 Layer Rendering Order}

\begin{lstlisting}[language=Python]
1. Ocean fill (biome-ocean, flat)
2. Elevation gradient (if elevation layer active)
3. Biome colors (if biome layer active) — per-cell fill
4. River overlay: blue strokes, width=1, opacity proportional to flow
5. Lake overlay: semi-transparent biome-lake fill
6. Resource heatmap (if resource layer active) — color mapped per resource
7. Habitability heatmap (if habitability layer active) — opacity 0.5
8. Death density heatmap (if deaths layer active) — opacity 0.35
9. Colony territory outlines (thin dashed, colony color, convex hull)
10. Agent dots
11. Hovered cell highlight (1px white border around cell)
12. Selected agent ring
\end{lstlisting}

\subsubsection*{9.3 Hover Tooltip}

\begin{lstlisting}[language=Python]
Position: follows cursor (offset: +12px right, +8px below)
Overflow: flip to left if near right edge

Content:
  [Biome name] — hsl(biome_color)
  Elevation: 0.62
  Rainfall:  823 mm/yr
  Temperature: 12.4°C
  Wood: ████░░ 0.68
  Stone: ██░░░░ 0.31
  Iron: ░░░░░░ 0.02
  [Causal explainer text if available]
\end{lstlisting}

\noindent\rule{\textwidth}{0.4pt}

\textbackslash{}textit{End of GDS v1.0}

\noindent\rule{\textwidth}{0.4pt}
\textbackslash{}textbf{Document produced by Antigravity AI, July 2026}
\textbackslash{}textbf{Companion: GPS  GES  GCR}

\end{document}
